% Generated draft: requires mathematical and source review.
\documentclass[11pt]{article}
\usepackage[margin=25mm]{geometry}
\usepackage{fontspec}
\setmainfont{Latin Modern Roman}
\usepackage{amsmath,amssymb,mathtools,unicode-math}
\setmathfont{Latin Modern Math}
\usepackage{enumitem,needspace,xurl,hyperref}
\setlength{\parindent}{0pt}
\setlength{\parskip}{5pt}
\setlength{\emergencystretch}{3em}
\newcommand{\N}{\mathbb{N}}
\newcommand{\Z}{\mathbb{Z}}
\newcommand{\Q}{\mathbb{Q}}
\newcommand{\R}{\mathbb{R}}
\newcommand{\C}{\mathbb{C}}
\newcommand{\F}{\mathbb{F}}
\newcommand{\A}{\mathbb{A}}
\newcommand{\PP}{\mathbb{P}}
\newcommand{\E}{\mathbb{E}}
\newcommand{\eps}{\varepsilon}
\newcommand{\dd}{\,\mathrm{d}}
\begin{document}
\section*{UnsolvedMath ID: 30003966}
\textbf{Source problem code:} catalogue entry
\subsection{Definitions and mathematical statement}
For the conventional setting of this question, let \(\Sigma_g\) be a closed, connected, oriented surface of genus \(g\geq 2\), and let \(n\geq 2\) be an integer. These surface hypotheses and parameter ranges are not explicitly supplied in the source; the formulation below is conditional on these conventions.

Let \(T(\Sigma_g)\) be the Teichm\"uller space of marked complex structures on \(\Sigma_g\), and let \(\operatorname{Mod}(\Sigma_g)\) be the group of isotopy classes of orientation-preserving diffeomorphisms of \(\Sigma_g\). Write \(\operatorname{Hit}_n(\Sigma_g)\) for a Hitchin component of the \(\operatorname{PSL}(n,\R)\) character variety of \(\pi_1(\Sigma_g)\): a connected component containing a representation obtained by composing a discrete faithful Fuchsian representation into \(\operatorname{PSL}(2,\R)\) with the principal homomorphism into \(\operatorname{PSL}(n,\R)\). Equip Teichm\"uller space and the Hitchin component with their natural mapping class group actions induced by changing markings.

For a marked Riemann surface \(X\in T(\Sigma_g)\), let \(K_X\) be its holomorphic cotangent line bundle, let \(K_X^i=K_X^{\otimes i}\), and let \(H^0(X,K_X^i)\) be the complex vector space of holomorphic sections of \(K_X^i\). Set
\[
V_X=\bigoplus_{i=3}^{n}H^0(X,K_X^i).
\]
For \(n=2\), this is the zero vector space by the empty direct sum convention.

The supplied question asks whether there exists a projection
\[
p:\operatorname{Hit}_n(\Sigma_g)\longrightarrow T(\Sigma_g)
\]
that is surjective and mapping-class-group-equivariant, meaning
\[
p(\gamma\cdot\rho)=\gamma\cdot p(\rho)
\qquad
\text{for every }\gamma\in\operatorname{Mod}(\Sigma_g)
\text{ and }\rho\in\operatorname{Hit}_n(\Sigma_g),
\]
and whose fibers satisfy
\[
p^{-1}(X)\cong V_X
\qquad\text{for every }X\in T(\Sigma_g).
\]
Here \(\cong\) records the requested fiber identification without resolving its category. The source does not specify whether the projection must be continuous, smooth, locally trivial, or satisfy another regularity condition; nor does it specify whether the fiber identifications are homeomorphisms, diffeomorphisms, or identifications with additional structure. It also does not specify compatibility of those identifications as \(X\) varies. Thus the displayed conditions preserve the mathematical target, but a fully precise existence statement requires these qualifications to be settled. No section, retraction property, or vector bundle structure is asserted.
\subsection{Short English statement}
The question asks whether the Hitchin component admits a mapping-class-group-equivariant projection to Teichm\"uller space such that the fiber over every marked Riemann surface \(X\) is identified with the space of holomorphic differentials of degrees \(3,\ldots,n\) on \(X\). The surface hypotheses, parameter quantifiers, regularity of the projection, and meaning of the fiber identification require clarification. The supplied registry labels the problem open; this draft does not independently verify that status.
\subsection{Sources}
\begin{itemize}
\item Local UnsolvedMath registry: supplied complete record for permanent ID \(30003966\), entitled \emph{Teichm\"uller Projections of Hitchin Components}, in the Oberwolfach Reports Open Problems collection; exact source problem code \texttt{catalogue entry}. The record gives proposed year 2018 and names no individual proposer.
\item Supplied report citation: \emph{Flat Surfaces and Algebraic Curves (2019)}, Oberwolfach Reports, pp.\ 2583--2650, DOI \texttt{10.4171/owr/2018/42}. Supplied source and literature link: \url{https://doi.org/10.4171/owr/2018/42}. The supplied background describes Oberwolfach Reports as documenting mathematical workshops and their open-problem discussions.
\item Supplied literature reference, labeled \emph{Unstable minimal surfaces}: \url{https://annals.math.princeton.edu/articles/22504}.
\item Supplied literature reference, labeled \emph{Hitchin component fibration}: \url{https://arxiv.org/abs/2208.04885}.
\item Supplied literature assessment, checked on \texttt{2026-08-21}: the record labels the problem open and states that rank two cases are proved, that later work disproves a related minimal map uniqueness conjecture in higher rank, and that the exact equivariant fibration and fiber statement remains unresolved. This is supplied dated curation, not an independently verified literature assessment; its time-sensitive status and cited claims require independent checking. The record also labels the statement \texttt{corrected\_verified} and reports a manual semantic, typographic, and LaTeX review; these are supplied editorial assertions, not mathematical verification by this draft.
\end{itemize}
\end{document}
